\documentclass[../main.tex]{subfiles}

\begin{document}

\section{Diseño del sistema}

\subsection{Requisitos del sistema}

En esta subsección se definen los requisitos que guían el diseño del sistema de sonificación EEG--MIDI. El propósito no es describir todavía la solución implementada a nivel de código, sino delimitar qué capacidades debe ofrecer el prototipo y qué condiciones generales debe respetar la arquitectura. La concreción de estos requisitos en módulos, contratos, parámetros y funciones se desarrolla posteriormente en el capítulo de Implementación.

El sistema se plantea como un prototipo experimental y no clínico. Su objetivo no es realizar diagnóstico médico ni interpretar patologías neurológicas, sino adquirir señales EEG reales, procesarlas durante la ejecución y utilizarlas como entrada para una generación musical enviada mediante una salida MIDI física.

Desde el punto de vista de diseño, se busca una arquitectura modular. La adquisición, el transporte de datos, el procesamiento digital, la evaluación de calidad, la sonificación, la salida MIDI y la visualización auxiliar se consideran bloques diferenciados. Esta separación permite razonar sobre la función de cada etapa sin introducir todavía los detalles internos de la implementación.

\subsubsection{Requisitos funcionales}

Los requisitos funcionales definen las capacidades que debe ofrecer el sistema a lo largo de la cadena EEG--MIDI.

\begin{description}[leftmargin=!,labelwidth=2cm]

    \item[RF-1] \textbf{Adquisición EEG en vivo.} El sistema debe adquirir señales EEG reales procedentes de un usuario durante una sesión de funcionamiento.

    La señal utilizada por la sonificación debe proceder de una captura en vivo y no de la reproducción \textit{offline} de registros previamente almacenados. Este requisito responde al objetivo principal del proyecto: transformar actividad EEG real en eventos musicales durante la ejecución del prototipo.

    \item[RF-2] \textbf{Adquisición multicanal ampliable.} El sistema debe contemplar una adquisición estructurada en varios canales, aunque una configuración concreta de validación pueda utilizar solo un canal principal.

    Este requisito permite que la arquitectura no quede limitada a un único montaje experimental y facilita futuras ampliaciones del sistema.

    \item[RF-3] \textbf{Coherencia temporal de la señal.} El sistema debe conservar una referencia temporal coherente entre adquisición, transmisión, procesamiento y validación.

    La frecuencia de muestreo, la organización temporal de las muestras y las ventanas de análisis deben mantenerse de forma consistente en toda la cadena.

    \item[RF-4] \textbf{Compatibilidad con electrodos de superficie.} El sistema debe permitir el uso de electrodos de superficie en un entorno experimental controlado.

    El diseño debe admitir configuraciones compatibles con la adquisición EEG, dejando el montaje físico concreto como una decisión experimental desarrollada en la implementación y la validación.

    \item[RF-5] \textbf{Transporte continuo de datos.} El sistema debe transportar de forma continua las muestras digitalizadas desde el bloque de adquisición hacia el entorno de procesamiento.

    El transporte debe permitir mantener el flujo de señal durante la sesión y detectar posibles pérdidas, discontinuidades o datos incoherentes.

    \item[RF-6] \textbf{Contrato estable entre bloques.} El sistema debe definir formatos de intercambio estables entre adquisición, procesamiento, sonificación, visualización y salida MIDI.

    Estos contratos deben permitir que cada bloque consuma la información necesaria sin depender de los detalles internos de los demás. Los nombres y formatos concretos de dichos contratos se detallan en el capítulo de Implementación.

    \item[RF-7] \textbf{Procesamiento por ventanas temporales.} El sistema debe procesar la señal EEG a partir de ventanas temporales recientes.

    Esta decisión evita basar la sonificación en muestras aisladas y permite obtener una representación más estable de la actividad registrada.

    \item[RF-8] \textbf{Extracción de características EEG.} El sistema debe extraer características temporales y espectrales de la señal adquirida.

    Estas características deben resumir la información relevante de la señal para que pueda utilizarse como entrada de la sonificación, sin que el motor musical dependa directamente de la forma de onda cruda.

    \item[RF-9] \textbf{Análisis por bandas EEG.} El sistema debe obtener información asociada a bandas EEG de interés.

    El análisis por bandas permite describir la distribución espectral de la señal y utilizarla como base para modular parámetros musicales de alto nivel. La definición concreta de las bandas y de las variables calculadas se desarrolla en la implementación.

    \item[RF-10] \textbf{Evaluación de calidad de señal.} El sistema debe valorar la fiabilidad de la señal antes de utilizarla en la sonificación.

    Esta evaluación debe permitir distinguir ventanas utilizables de ventanas afectadas por ruido, artefactos, saturación, mala conexión o problemas de transporte.

    \item[RF-11] \textbf{Sonificación condicionada por la calidad.} El sistema debe evitar que una ventana de señal no fiable tenga la misma influencia musical que una ventana válida.

    La calidad de la señal debe poder dejar pasar, atenuar o bloquear la influencia de las características EEG sobre la generación musical. El mecanismo concreto empleado para ello se describe en la implementación.

    \item[RF-12] \textbf{Generación musical a partir de características EEG.} El sistema debe transformar las características extraídas de la señal en parámetros musicales.

    Estos parámetros no deben interpretarse como variables clínicas, sino como controles musicales que modifican aspectos como actividad, densidad, registro, estabilidad, dinámica o probabilidad de generación de notas.

    \item[RF-13] \textbf{Organización musical estructurada.} El sistema debe generar música mediante una estructura organizada y no mediante una conversión directa de muestra a nota.

    La señal EEG debe influir en decisiones musicales de alto nivel que posteriormente se concreten en compases, notas y eventos MIDI. La estructura concreta de generación se desarrolla en el capítulo de Implementación.

    \item[RF-14] \textbf{Salida MIDI física.} El sistema debe emitir la salida musical mediante una interfaz MIDI física.

    La salida principal del prototipo debe ser una secuencia real de mensajes MIDI enviada hacia un dispositivo externo compatible, no únicamente una visualización en pantalla.

    \item[RF-15] \textbf{Eventos MIDI completos y seguros.} El sistema debe generar eventos MIDI con inicio y finalización definidos.

    Esta condición evita notas incompletas o sostenidas indefinidamente y permite incluir acciones de seguridad para liberar la salida musical cuando sea necesario.

    \item[RF-16] \textbf{Monitorización auxiliar.} El sistema debe disponer de una interfaz auxiliar para observar el estado del prototipo.

    Esta interfaz debe facilitar la supervisión de adquisición, procesamiento, calidad, sonificación y salida MIDI, sin sustituir a la ruta principal EEG--MIDI ni duplicar los cálculos críticos.

    \item[RF-17] \textbf{Registro y trazabilidad.} El sistema debe permitir registrar información útil para validación, depuración y análisis posterior.

    Las capturas, métricas y estados publicados deben ayudar a reconstruir el comportamiento del sistema durante las pruebas.

\end{description}

\subsubsection{Requisitos no funcionales}

Los requisitos no funcionales recogen las condiciones de calidad, seguridad, operación y mantenimiento que debe cumplir el prototipo durante su uso experimental.

\begin{description}[leftmargin=!,labelwidth=2.2cm]

    \item[RNF-1] \textbf{Funcionamiento en tiempo real experimental.} La cadena de adquisición, procesamiento, sonificación y salida MIDI debe operar de forma continua durante la sesión.

    El sistema no busca certificación clínica ni garantías propias de un dispositivo médico, pero sí debe disponer de margen temporal suficiente para funcionar de forma fluida en el contexto del prototipo.

    \item[RNF-2] \textbf{Uso no clínico.} El sistema debe orientarse a un uso experimental, educativo y de investigación no clínica.

    No debe presentarse como dispositivo médico ni como herramienta diagnóstica. La interpretación de la señal se limita al contexto de la sonificación y la validación técnica del prototipo.

    \item[RNF-3] \textbf{Seguridad básica del montaje.} El diseño debe considerar las precauciones necesarias para trabajar con señales biopotenciales en un entorno controlado.

    Las decisiones de conexión de electrodos, alimentación y uso del prototipo deben plantearse desde una perspectiva de baja tensión, aislamiento razonable y uso experimental supervisado.

    \item[RNF-4] \textbf{Robustez frente al ruido y artefactos.} El diseño debe tener en cuenta la baja amplitud de la señal EEG y su sensibilidad a interferencias eléctricas, movimiento, actividad muscular y problemas de contacto electrodo--piel.

    Esta robustez debe abordarse tanto desde la adquisición como desde la evaluación de calidad y la monitorización del sistema.

    \item[RNF-5] \textbf{Modularidad.} El sistema debe separarse en bloques funcionales diferenciados y con responsabilidades claras.

    Esta modularidad facilita el análisis de cada etapa y permite evolucionar adquisición, procesamiento, sonificación, MIDI o visualización sin rediseñar toda la arquitectura.

    \item[RNF-6] \textbf{Escalabilidad.} La arquitectura debe permitir futuras ampliaciones.

    Entre las extensiones posibles se incluyen nuevos canales EEG, otros montajes de electrodos, nuevas visualizaciones, nuevos mapeos musicales o mejoras en la evaluación de calidad.

    \item[RNF-7] \textbf{Bajo coste y accesibilidad.} El prototipo debe priorizar componentes y herramientas accesibles en un contexto académico y de prototipado.

    Esta condición afecta tanto a la selección del hardware como al uso de herramientas de desarrollo y validación reproducibles con recursos razonables.

    \item[RNF-8] \textbf{Replicabilidad y documentación.} El sistema debe documentarse de forma que pueda ser comprendido, reproducido, validado y modificado posteriormente.

    La memoria, el código, los diagramas, las capturas y las herramientas de validación deben permitir reconstruir el flujo seguido desde la adquisición EEG hasta la salida MIDI.

    \item[RNF-9] \textbf{Mantenibilidad.} El código y la arquitectura deben organizarse de forma clara, reduciendo dependencias innecesarias entre módulos.

    El sistema debe facilitar ajustes posteriores en el procesamiento, la sonificación, la interfaz auxiliar o las herramientas de validación sin comprometer la ruta principal EEG--MIDI.

    \item[RNF-10] \textbf{Observabilidad.} El sistema debe exponer suficiente información interna para facilitar depuración, validación y análisis posterior.

    La visualización auxiliar, las capturas, los estados publicados y los registros de rendimiento permiten observar el comportamiento del sistema sin alterar la ruta principal de adquisición, procesamiento y salida MIDI.

    \item[RNF-11] \textbf{Portabilidad del prototipo.} El sistema debe poder utilizarse como prototipo compacto y transportable en un entorno experimental controlado.

    Esta portabilidad se refiere tanto al montaje físico como a la ejecución integrada de las capas embebida y de procesamiento.

\end{description}

\subsubsection{Decisiones de diseño derivadas}

A partir de los requisitos anteriores se adoptan varias decisiones de diseño que condicionan el resto de la arquitectura. En primer lugar, se separa la adquisición embebida del procesamiento de alto nivel. La capa cercana al hardware se encarga de obtener y preparar la señal, mientras que la capa de procesamiento trabaja con datos ya estructurados.

En segundo lugar, el sistema se organiza mediante contratos de intercambio entre bloques. El diseño no exige que cada etapa conozca la implementación interna de las demás, sino que reciba y entregue información con una estructura definida. Los contratos concretos y sus campos se detallan en el capítulo de Implementación.

En tercer lugar, el procesamiento se basa en ventanas temporales de señal y no en muestras aisladas. Esta elección permite obtener características más estables y evita una traducción excesivamente directa entre señal instantánea y salida musical.

En cuarto lugar, la calidad de la señal se incorpora como criterio de diseño. La sonificación debe depender no solo de las características extraídas, sino también de la confianza que el sistema tenga en la ventana analizada.

En quinto lugar, la generación musical se concibe como una estructura jerárquica. Las características EEG se transforman primero en controles musicales y después en elementos musicales concretos. Esta organización favorece una salida más coherente que un mapeo directo entre bandas EEG y notas aisladas.

Por último, la salida principal se define como MIDI físico. La interfaz web, la visualización de notas, las capturas y las métricas se consideran herramientas auxiliares de observación, validación y depuración. La concreción de estas decisiones, incluyendo módulos, parámetros temporales, contratos y funciones, se desarrolla en el capítulo de Implementación.

\subsection{Arquitectura general del sistema}

% Pendiente de redacción propia.
% Describir el flujo global: electrodos -> ADS1299 -> Arduino UNO Q -> backend Python -> sonificación -> salida MIDI / interfaz web.

\subsection{Bloque de adquisición EEG}

% Pendiente de redacción propia.
% Describir electrodos, ADS1299-4PAG, canales, referencia, BIAS/RLD, ganancia, frecuencia de muestreo y lectura sincronizada.

\subsection{Bloque de comunicación microcontrolador--Linux}

% Pendiente de redacción propia.
% Describir el envío de bloques desde el firmware hacia el entorno Python en la Arduino UNO Q.

\subsection{Bloque de procesamiento digital}

% Pendiente de redacción propia.
% Describir filtrado, ventanas, buffer temporal, PSD, potencia por bandas, picos espectrales y métricas de calidad.

\subsection{Bloque de sonificación EEG--MIDI}

% Pendiente de redacción propia.
% Describir segmentación EEG, construcción de segmentos musicales, generación de compases, notas y eventos MIDI.

\subsection{Bloque de visualización e interacción}

% Pendiente de redacción propia.
% Describir interfaz web HTML, métricas mostradas, piano scroll, selección de canal activo y visualización en matriz LED.

\subsection{Diseño electrónico y PCB}

% Pendiente de redacción propia.
% Describir conexiones principales: Arduino UNO Q, ADS1299, electrodos, salida MIDI, alimentación y consideraciones de seguridad.

\end{document}